\documentclass[a4paper]{article}

\usepackage[paperwidth=100mm,margin=.1mm,ignoreall]{geometry} % prose
%\usepackage[paperwidth=62.5mm,margin=.1mm,ignoreall]{geometry}  % poem
\usepackage[dvipsnames,svgnames]{xcolor}
\usepackage{graphicx,listings}

\usepackage[path={./fonts/}]{rit-fonts}
\newfontfamily\en[Path={./fonts/},Extension=.ttf]%
  {Linux-Libertine}
  [BoldFont=*-Bold]
\def\thesection{\en\arabic{section}}
\def\thesubsection{\en\thesection.\arabic{subsection}}
\def\f#1{\gdef\thef{\en #1}}
\newfontfamily\lstfont[Scale=.8]{MPLUSCodeLatin}
\usepackage{tcolorbox}

\colorlet{bgcol}{white}
\colorlet{frcol}{black!10}
\def\thetitle{Regular}

\makeatletter
\def\title#1{\def\next{#1}\xdef\thetitle
  {\expandafter\strip@prefix\meaning\next}}
\makeatother

\newtcolorbox{prosebox}{colback=bgcol,width=99.8mm, %.63\linewidth,
    halign title=flush left,
  	before=, after=\par\newpage, 
 % 	before skip=0\baselineskip, after skip=0\baselineskip,
  	colbacktitle={black!10},
  	coltitle=black, 
  	fonttitle=\lstfont\footnotesize, 
  	box align=top,
  	colframe=frcol,
  	title={\thetitle}}

%\addfontfeatures{Scale=2,Color=red,Numbers=Lining}

\def\myprose{പ്രതിരോധരാഷ്ട്രീയവും ആധുനിക കലാപാരമ്പര്യങ്ങളിലെ 
  വിനിമയവൈഭവങ്ങളും തമ്മിലിണക്കി, അവരവരുടെ ദേശീയ കലാസാഹിത്യത്തിന്റെ 
  പരമ്പരാഗത രൂപങ്ങളെ ഭാവുകത്വപരമായി ഉയർത്തി, നീതിയോടു് ആഭിമുഖ്യം വളർത്തി, 
  ഇങ്ങനെ ഒരു പ്രതിസംസ്കാരത്തിന്റെ ജനകീയവും അഗാധവും സമകാലികവും 
  ഭൗതികവുമായ ജാഗ്രത പാകി.}

\def\halfprose{പ്രതിരോധരാഷ്ട്രീയവും ആധുനിക കലാപാരമ്പര്യങ്ങളിലെ 
  വിനിമയവൈഭവങ്ങളും തമ്മിൽ  \dots}

\usepackage[active,luatex,tightpage]{preview}
\renewcommand\PreviewBbAdjust{0pt -5pt 0pt 0pt}
\PreviewEnvironment{prosebox}

  
\begin{document}
\ifluatex\prehyphenchar="0000\else\defaulthyphenchar="0000\fi
\pagestyle{empty}

\parindent0pt
\parskip=1\baselineskip plus 2pt minus 2pt

\hyphenpenalty=0
\lefthyphenmin=3
\righthyphenmin=3
\tolerance=500
\emergencystretch=1em

\title{RIT Rachana Regular}
\begin{prosebox}
\Rachana
\myprose
\end{prosebox}

\title{RIT Panmana Regular}
\begin{prosebox}
\Panmana
\myprose
\end{prosebox}

\title{RIT MeeraNew}
\begin{prosebox}
\Meera
\myprose
\end{prosebox}

\title{RIT TN Joy}
\begin{prosebox}
\Joy
\myprose
\end{prosebox}

\title{\usepackage[ScaleRM=1.5,...]{rit-fonts}}
\begin{prosebox}
  \Rachana
  \addfontfeatures{Scale=1.5}
  \advance\baselineskip 5pt
  \ifluatex\prehyphenchar="0000\else\defaulthyphenchar="0000\fi
\halfprose
\end{prosebox}


\title{\usepackage[Numbers=OldStyle,...]{rit-fonts}}
\begin{prosebox}
  \Rachana \addfontfeatures{Numbers=OldStyle}
  \ifluatex\prehyphenchar="0000\else\defaulthyphenchar="0000\fi രചനയിൽ
  പഴയരീതിയിലുള്ള അക്കങ്ങളുണ്ടു്. 1234567890. ഇവ സാധാ\-രണ അക്കങ്ങളെ അപേക്ഷിച്ചു്
  പാഠവുമായി ഒത്തുപോകുന്നതാണു്.
\end{prosebox}


\end{document}

